\section{Dogma, Doctrine, Typology, and Title}

The four-dogma corpus is not an exhaustive Mariology. Catholic teaching includes other authoritative doctrines; liturgy and the Fathers supply typologies; popular piety uses titles; theologians propose conclusions. Precision about category protects rather than diminishes these truths. Calling every honored expression a dogma would erase the Church's own distinctions and make responsible development impossible.

\begin{fourstable}{Teaching or title}{Proper classification}{Secure content}{Boundary}
Lifelong personal holiness & Authoritative doctrine closely joined to the Immaculate Conception & By singular grace Mary remained free from every personal sin. & Not identical to the 1854 definition, whose exact object is original sin at the first instant.\\
Mother of the Church & Authoritative ecclesial title and doctrine & Mother of Christ the Head, she is maternally related in grace to his members. & She does not generate the Church's divine life independently or replace the Spirit.\\
Spiritual motherhood & Conciliar and ordinary magisterial teaching & Her faith, obedience, charity, and intercession serve the birth and growth of believers. & Analogical to physical maternity and wholly dependent on Christ.\\
Queenship & Authoritative papal and liturgical doctrine & She shares, in a limited and analogous created manner, the royal glory of her Son. & Not equal sovereignty, natural omnipotence, or a separately defined coronation event.\\
New Eve & Patristic theological typology & Her obedient faith answers Eve's disobedience within Christ's recapitulation. & Christ alone is the saving new Adam; the parallel does not create two redeemers.\\
Ark / Daughter Zion & Biblical--patristic typologies & She bears the incarnate presence and personally receives Israel's promised King. & Not the exhaustive literal sense of Exodus, 2 Samuel, Zephaniah, or Revelation.\\
All-Holy (\emph{Panagia}) & Ancient traditional and liturgical title & Celebrates the sanctifying grace manifested throughout her vocation. & Older usage does not automatically state the technical first-instant definition.\\
Death or Dormition & Weighty and ancient theological-liturgical tradition & Mary is commonly held to have died in conformity with Christ before glorification. & Her death was deliberately not defined in 1950.\\
\end{fourstable}

\subsection{Lifelong freedom from personal sin}

The Immaculate Conception directly concerns original sin. Catholic doctrine also teaches that Mary, sustained by grace, remained free of every personal sin. The Council of Trent says no justified person can avoid all venial sins throughout life without a special privilege from God, ``as the Church holds concerning the Blessed Virgin'' (session VI, canon 23). Vatican II calls her ``all holy and free from every stain of sin'' and presents her fiat as an undivided embrace of God's will (LG 56). CCC 493 summarizes the Eastern and Western confession.

This does not mean Mary had no need to deliberate, grow in understanding, exercise faith, or suffer obscurity. Luke says she pondered what she did not yet comprehend fully (Luke 2:19, 50--51). Sinlessness is not omniscience. Nor should difficult Gospel scenes be solved by pretending they contain no human surprise. The Church's conclusion is that no personally culpable disorder ruptured her communion with God; it is not that her earthly consciousness was the beatific vision.

Historical honesty again matters. Certain Fathers interpreted the search for Jesus, Cana, or the Cross in ways that attributed hesitation or correction to Mary. Those local exegeses do not overturn the Church's mature judgment, but neither may they be rewritten as though they already asserted it. The doctrine is received through the whole Tradition and Magisterium, not the isolated opinion of every Father.

\subsection{Mother of the Church and spiritual motherhood}

At the close of the third session of Vatican II in 1964, Paul VI proclaimed Mary ``Mother of the Church,'' meaning mother of the whole Christian people, pastors and faithful. The title was later given universal liturgical expression through the memorial of the Blessed Virgin Mary, Mother of the Church. Its theological roots lie in her maternity of Christ, her association with his saving mission, John 19:25--27, and her prayer within the apostolic community (Acts 1:14).

LG 61 says that by conceiving Christ, bringing him forth, nourishing him, presenting him to the Father, and suffering with him as he died, Mary cooperated in the Savior's work through obedience, faith, hope, and charity, so that supernatural life might be restored to souls; for this reason she is mother to us in the order of grace. The phrase ``in the order of grace'' is decisive. God alone gives divine life as first cause. Mary's maternity is personal, moral, intercessory, exemplary, and dispositive within Christ's one mediation. She is not mother of the Holy Spirit, mother of the divine essence, or efficient source of adoption.

The Church herself is mother through preaching and sacrament. Mary does not replace this ecclesial maternity. She is its exemplary realization and maternal member. Nor does the Church replace Mary with a symbol: the personal relation of the mother of Jesus to his disciples grounds the typology.

\subsection{Queenship}

Mary's Queenship derives from Christ's kingship and her relation to him. Patristic and liturgical language calls the Mother of the King Queen; iconography places her beside the enthroned Son. Pius XII's \emph{Ad caeli Reginam} (1954) synthesizes this tradition and expressly calls Mary's royal dignity a limited and analogous participation in Christ's kingship (especially 31, 34, 39). Vatican II says she was exalted by the Lord as Queen of all so that she might be more fully conformed to her Son (LG 59); CCC 966 receives the teaching.

Derived Queenship is neither constitutional competition nor absolute sovereignty. Mary does not exercise a jurisdiction independent of Christ, issue grace from her own treasury, or possess omnipotence by nature. Expressions such as ``suppliant omnipotence'' can mean that loving intercession is powerful by God's generosity; taken literally as a second omnipotence, they are false.

The Glorious Mysteries often name a ``Coronation of Mary.'' That devotional mystery contemplates the doctrinal reality of her Queenship and glory. Scripture and dogma do not narrate a separate visible ceremony after the Assumption. Psalm 45 and Revelation 12 provide typological imagery; they should not be converted into an additional historical dogma.

\subsection{New Eve, Ark, Daughter Zion, and all-holy}

Typology recognizes correspondences authored within the one divine economy. New Eve concerns obedient faith within Christ's recapitulation. Ark concerns the bearer of divine presence. Daughter Zion concerns Israel's faithful reception of the King. The woman of Revelation 12 concerns the people of God and, within that corporate symbol, the personal mother of the Messiah. No one image exhausts Mary, and none can be detached from Christ and the Church.

The Eastern title \emph{Panagia}, All-Holy, witnesses to the Church's praise of Mary's sanctification. Its broad antiquity supports the doctrinal matrix of unique holiness. Historical discipline asks what a particular author or liturgical text means by purity, stain, holiness, or consecration. A word that resembles ``immaculate'' in translation is not automatically a claim about original sin at the first instant.

\subsection{Private revelation and later titles}

Reported apparitions may use titles or images that intensify reception of established teaching. Even when competent authority permits or encourages devotion, private revelation does not complete the apostolic deposit, produce a new dogma, or supply the dogma's formal proof (DV 4; CCC 66--67). The 1858 Lourdes reports followed the 1854 definition; they did not confer its authority. Devotional promises and alleged messages must be judged according to their own ecclesial status and cannot silently govern universal Mariology.

The same rule applies to proposed dogmatic titles. Long use, presence in a papal address, an approved office, saintly advocacy, theological petitions, or popular signatures can be historically and theologically significant. None by itself is a solemn definition. The Church's current doctrinal judgment controls present constructive use, especially for ``Co-redemptrix'' and ``Mediatrix of All Graces.''

\subsection{Why classification serves devotion}

Category discipline is not minimalism. A typology can illuminate Scripture without being its literal sense. An authoritative doctrine can bind Catholic understanding without having a separate ex cathedra formula. A liturgical title can richly express worship without answering every scholastic question. A theological opinion can be venerable without being obligatory. Such distinctions let every truth bear its proper weight and prevent devotion from being built on claims the Church has not made.
